\chapter{Evaluating Supervised Models}
\label{ch:evaluation}

\chapterguide%
  {\chapref{ch:data} and \chapref{ch:statistics}: honest holdouts, label types, averages, spread, and sampling uncertainty.}
  {choose metrics, thresholds, and validation designs; assess probabilities, uncertainty, subgroups, and shift.}

Evaluate with realistic splits, relevant baselines, and metrics tied to error costs.

\section{Evaluation begins with the real objective}

Choose metrics for the supported decision and its error costs.

A useful evaluation plan states:

\begin{itemize}
  \item the unit being predicted and the population of interest;
  \item how the data split imitates future use;
  \item the baseline that must be beaten;
  \item the primary metric used for model selection;
  \item secondary metrics and subgroup checks;
  \item acceptable failure rates and operational constraints.
\end{itemize}

Prespecify the primary metric to avoid selective reporting.

\section{Regression metrics}
\label{sec:regression-metrics}

For true values $y_i$ and predictions $\hat{y}_i$, mean squared error and root mean squared error are
\[
  \operatorname{MSE}
  =\frac{1}{n}\sum_{i=1}^{n}{(y_i-\hat{y}_i)}^2,
  \qquad
  \operatorname{RMSE}=\sqrt{\operatorname{MSE}}.
\]
MSE strongly penalizes large errors. RMSE is easier to interpret because it uses the same units as the label.

\begin{workedexample}
Five delivery times, in minutes. Four predictions are nearly right; one is badly wrong.

\begin{mltable}\small
\begin{tabular}{lccccc}
\toprule
\rowcolor{MLPaleBlue!92}
True $y$            & 10 & 12 & 14 & 16 & 48 \\
Predicted $\hat{y}$ & 11 & 13 & 13 & 15 & 20 \\
Error $y-\hat{y}$   & $-1$ & $-1$ & $+1$ & $+1$ & $+28$ \\
\bottomrule
\end{tabular}
\end{mltable}

The absolute errors are $1,1,1,1,28$, so
\[
  \operatorname{MAE}=\frac{1+1+1+1+28}{5}=\frac{32}{5}=6.4,
  \qquad
  \operatorname{MedAE}=1.
\]
Squaring first gives $1,1,1,1,784$, so
\[
  \operatorname{MSE}=\frac{788}{5}=157.6,
  \qquad
  \operatorname{RMSE}=\sqrt{157.6}\approx 12.55.
\]
For $R^2$, the mean of the true values is $\bar{y}=100/5=20$, and
\[
  \sum{(y_i-\bar{y})}^2=100+64+36+16+784=1000,
  \qquad
  R^2=1-\frac{788}{1000}=0.212.
\]

MedAE $=1$, MAE $=6.4$, and RMSE $\approx12.55$ describe different aspects of error. One row contributes $784/788=99.5\%$ of squared error; inspect extremes when RMSE greatly exceeds MAE.
\end{workedexample}

\begin{pitfall}
Binning a continuous label changes the task to classification. Do so only when predefined categories match the decision.
\end{pitfall}

\section{Additional regression metrics}

Mean absolute error is
\[
  \operatorname{MAE}
  =\frac{1}{n}\sum_{i=1}^{n}|y_i-\hat{y}_i|.
\]
MAE weights errors linearly; median absolute error describes a typical miss and is more robust.

Mean absolute percentage error is
\[
  \operatorname{MAPE}
  =\frac{100}{n}\sum_{i=1}^{n}
  \left|\frac{y_i-\hat{y}_i}{y_i}\right|.
\]
MAPE is undefined at zero, unstable near zero, and asymmetric between over- and underprediction.

For asymmetric real-world costs, a custom loss or quantile loss may be more appropriate. Quantile regression can estimate an interval of plausible outcomes rather than only a conditional mean.

\section{Classification metrics: start with the confusion matrix}
\label{sec:classification-metrics}

Almost every classification metric is a different way of reading the same four boxes.

\begin{mlfigure}
\begin{tikzpicture}[font=\scriptsize]
% Axis titles and class labels
\node[font=\small\bfseries\color{MLNavy}] at (5.225,4.95) {Predicted class};
\node[font=\small\bfseries\color{MLNavy},rotate=90] at (1.00,2.50) {Actual class};
\node[draw=MLBlue!35,fill=MLPaleBlue,rounded corners=1mm,
      minimum width=3.10cm,minimum height=0.48cm,text=MLNavy]
  at (3.55,4.50) {Positive};
\node[draw=MLBlue!35,fill=MLPaleBlue,rounded corners=1mm,
      minimum width=3.10cm,minimum height=0.48cm,text=MLNavy]
  at (6.90,4.50) {Negative};
\node[draw=MLBlue!35,fill=MLPaleBlue,rounded corners=1mm,
      minimum width=1.50cm,minimum height=0.48cm,text=MLNavy,rotate=90]
  at (1.65,3.40) {Positive};
\node[draw=MLBlue!35,fill=MLPaleBlue,rounded corners=1mm,
      minimum width=1.50cm,minimum height=0.48cm,text=MLNavy,rotate=90]
  at (1.65,1.60) {Negative};

% Confusion-matrix cells
\filldraw[fill=MLPaleTeal,draw=MLGreen,rounded corners=1mm]
  (2.00,2.65) rectangle (5.10,4.15);
\filldraw[fill=MLPaleRed,draw=MLRed,rounded corners=1mm]
  (5.35,2.65) rectangle (8.45,4.15);
\filldraw[fill=MLPaleRed,draw=MLRed,rounded corners=1mm]
  (2.00,0.85) rectangle (5.10,2.35);
\filldraw[fill=MLPaleTeal,draw=MLGreen,rounded corners=1mm]
  (5.35,0.85) rectangle (8.45,2.35);

\node[align=center,text=MLNavy] at (3.55,3.40)
  {\textbf{True positive (TP)}\\caught it};
\node[align=center,text=MLNavy] at (6.90,3.40)
  {\textbf{False negative (FN)}\\missed it};
\node[align=center,text=MLNavy] at (3.55,1.60)
  {\textbf{False positive (FP)}\\false alarm};
\node[align=center,text=MLNavy] at (6.90,1.60)
  {\textbf{True negative (TN)}\\correctly ignored};

% Metric callouts connected to the entries they use
\node[draw=MLOrange,fill=MLPaleOrange,rounded corners=1mm,
      align=left,text width=3.25cm,minimum height=1.15cm,text=MLNavy]
  (recall) at (10.95,3.40)
  {\textbf{Recall} reads across the actual-positive row: $TP+FN$.};
\draw[MLOrange,thick] (8.67,2.65)--(8.67,4.15);
\draw[MLOrange,thick] (8.55,2.65)--(8.67,2.65)
                           (8.55,4.15)--(8.67,4.15);
\draw[-{Stealth[length=1.8mm]},MLOrange,thick]
  (8.67,3.40)--(recall.west);

\node[draw=MLBlue,fill=MLPaleBlue,rounded corners=1mm,
      align=left,text width=3.25cm,minimum height=1.15cm,text=MLNavy]
  (precision) at (10.95,1.45)
  {\textbf{Precision} reads down the predicted-positive column: $TP+FP$.};
\draw[MLBlue,thick] (2.00,0.63)--(5.10,0.63);
\draw[MLBlue,thick] (2.00,0.63)--(2.00,0.75)
                         (5.10,0.63)--(5.10,0.75);
\draw[-{Stealth[length=1.8mm]},MLBlue,thick]
  (3.55,0.63)--(3.55,0.42)--(8.85,0.42)--(8.85,1.45)--(precision.west);
\end{tikzpicture}
\end{mlfigure}

\begin{analogy}
Precision measures how many alerts are correct; recall measures how many positives are detected. Always alerting sacrifices precision; never alerting gives zero recall and undefined precision.
\end{analogy}

\begin{mltable}
\begin{tabularx}{\linewidth}{@{}>{\bfseries\leavevmode\color{MLNavy}}p{0.18\linewidth}p{0.26\linewidth}X@{}}
\toprule
\rowcolor{MLPaleBlue!92}
\normalfont\bfseries\leavevmode\color{MLNavy}Metric & \normalfont\bfseries\leavevmode\color{MLNavy}Formula & \normalfont\bfseries\leavevmode\color{MLNavy}Answers the question \\
\midrule
Accuracy & $\dfrac{TP+TN}{\text{all}}$ & How often is the model right overall? \\[0.5em]
Precision & $\dfrac{TP}{TP+FP}$ & When it says yes, is it right? \\[0.5em]
Recall & $\dfrac{TP}{TP+FN}$ & Does it find the ones that matter? \\[0.5em]
F1 & $2\cdot\dfrac{P\cdot R}{P+R}$ & One number balancing the two \\[0.5em]
Specificity & $\dfrac{TN}{TN+FP}$ & Does it leave the negatives alone? \\
\bottomrule
\end{tabularx}
\end{mltable}

\begin{pitfall}
Read confusion-matrix axis labels; actual/predicted conventions vary.
\end{pitfall}

\begin{workedexample}
Of 1,000 emails, 100 are spam. The filter flags 80, of which 60 are spam:
\[
  TP=60,\qquad FP=80-60=20,\qquad FN=100-60=40,\qquad TN=900-20=880.
\]
(Check the total: $60+20+40+880=1000$.) Now every metric is arithmetic:
\[
  \text{Accuracy}=\frac{60+880}{1000}=0.940,
  \qquad
  \text{Precision}=\frac{60}{80}=0.750,
  \qquad
  \text{Recall}=\frac{60}{100}=0.600,
\]
\[
  F_1=2\cdot\frac{0.750\times 0.600}{0.750+0.600}
     =\frac{0.900}{1.350}\approx 0.667.
\]

The no-alert baseline reaches 90\% accuracy; the model reaches 94\% but misses 40 spam emails. Recall exposes those misses.
\end{workedexample}

\section{Additional confusion-matrix measures}

In addition to precision and recall, useful classification measures include
\[
  \text{Specificity}=\frac{TN}{TN+FP},
  \qquad
  \text{False-positive rate}=\frac{FP}{FP+TN}.
\]
Balanced accuracy averages recall for the positive and negative classes:
\[
  \text{Balanced accuracy}
  =\frac{1}{2}
  (\text{Recall}+\text{Specificity}).
\]
It is more informative than ordinary accuracy when classes are imbalanced.

The Matthews correlation coefficient (MCC) summarizes the entire confusion matrix in one number:
\[
  \operatorname{MCC}
  =\frac{TP\cdot TN-FP\cdot FN}
  {\sqrt{(TP+FP)(TP+FN)(TN+FP)(TN+FN)}}.
\]
MCC ranges from $-1$ to $1$; near zero indicates little association. Using all four cells makes it useful under imbalance.

\begin{workedexample}
Continuing the spam filter ($TP=60$, $FP=20$, $FN=40$, $TN=880$):
\[
  \text{Specificity}=\frac{880}{900}\approx 0.978,
  \qquad
  \text{Balanced accuracy}=\frac{0.600+0.978}{2}\approx 0.789.
\]
Accuracy is $0.940$, dominated by 900 negatives. Balanced accuracy is $0.789$, weighting each class equally.

MCC uses all four cells at once. The numerator is
\[
  TP\cdot TN-FP\cdot FN=60(880)-20(40)=52{,}800-800=52{,}000,
\]
and the denominator is
\[
  \sqrt{(80)(100)(900)(920)}=\sqrt{6{,}624{,}000{,}000}\approx 81{,}388,
\]
so $\operatorname{MCC}\approx 52{,}000/81{,}388\approx 0.639$.

Accuracy, balanced accuracy, and MCC answer different questions; do not report accuracy alone under imbalance.
\end{workedexample}

Multiclass \term{macro averaging} weights classes equally; \term{micro averaging} pools decisions; \term{weighted averaging} weights classes by frequency.

\section{Thresholds, ROC curves, and precision-recall curves}

A probability model becomes a hard classifier only after a threshold is chosen. Lowering the threshold usually increases recall and false positives; raising it usually increases precision and false negatives.

The receiver operating characteristic (ROC) curve plots true-positive rate against false-positive rate across thresholds. ROC AUC is the probability that a randomly chosen positive example receives a higher score than a randomly chosen negative example.

Choose thresholds or batch size on validation predictions using costs or capacity. Lock the rule before test evaluation (Lab~2, Part~C).

% === VISUAL ADDITION: threshold-operating-point ===
\subsection{Thresholds move the operating point}

These six scores stay fixed; only the decision threshold changes.

\begin{mlfigure}
\begin{tikzpicture}[x=10.0cm,y=1cm,font=\scriptsize]
  \draw[-{Stealth[length=2mm]},MLMuted] (0,0) -- (1.02,0)
    node[right] {predicted probability};
  \foreach \x/\lab in {0/0,0.2/.2,0.4/.4,0.6/.6,0.8/.8,1/1}
    \draw[MLMuted] (\x,0.06)--(\x,-0.06) node[below=2pt] {\lab};

  % scored cases
  \fill[MLBlue] (0.18,0.48) circle (2.7pt);
  \draw[MLOrange,very thick] (0.32,0.48) circle (3.0pt);
  \fill[MLBlue] (0.41,0.48) circle (2.7pt);
  \draw[MLOrange,very thick] (0.46,0.48) circle (3.0pt);
  \fill[MLBlue] (0.64,0.48) circle (2.7pt);
  \draw[MLOrange,very thick] (0.84,0.48) circle (3.0pt);

  \node[anchor=west,text=MLBlue] at (0.03,1.10) {$\bullet$ actually negative};
  \node[anchor=west,text=MLOrange] at (0.34,1.10) {$\circ$ actually positive};

  \draw[MLNavy,very thick] (0.50,-0.28)--(0.50,0.92);
  \node[align=left,text=MLNavy,anchor=west] at (0.54,-0.48)
    {threshold $0.50$\\flags 2 cases};

  \draw[MLRed,very thick,dashed] (0.30,-0.48)--(0.30,0.92);
  \node[align=center,text=MLRed] at (0.15,-1.05)
    {lower to $0.30$\\catches more positives\\but also more negatives};

  \draw[-{Stealth[length=1.8mm]},MLRed,thick] (0.48,-0.18)--(0.32,-0.18);
\end{tikzpicture}
\end{mlfigure}

\begin{keyidea}
Lower thresholds increase positive decisions and cannot reduce recall. Higher thresholds reduce alerts but may miss positives; validate the cost/capacity trade-off.
\end{keyidea}
% === END VISUAL ADDITION ===

\subsection{Capacity and cost metrics}
\label{sec:capacity-metrics}

When only $k$ cases can be reviewed, rank all eligible cases and report
\[
\operatorname{precision@}k=\frac{\text{positives in top }k}{k},
\qquad
\operatorname{recall@}k=\frac{\text{positives in top }k}{\text{all positives}},
\]
\[
\operatorname{lift@}k=
\frac{\operatorname{precision@}k}{\text{positive prevalence}}.
\]
Report error costs when defensible. For fixed review capacity, evaluate the ranked batch, as in Lab~5's top 100 fraud alerts.

\begin{workedexample}
A team can call five of 20 customers. Ranked outcomes follow; $1$ denotes churn:
\begin{center}\ttfamily\small
1\quad 1\quad 0\quad 1\quad 0\ \ \vrule\ \ 0\quad 0\quad 1\quad 0\quad 0\quad 0\quad 0\quad 0\quad 0\quad 0\quad 0\quad 0\quad 0\quad 0\quad 0
\end{center}
The bar marks the capacity cutoff at $k=5$. There are $4$ churners in the batch of $20$, so the prevalence is $4/20=0.20$. Of the top five, three churned:
\[
  \operatorname{precision@}5=\frac{3}{5}=0.60,
  \qquad
  \operatorname{recall@}5=\frac{3}{4}=0.75,
  \qquad
  \operatorname{lift@}5=\frac{0.60}{0.20}=3.0.
\]
Randomly calling five customers would reach one churner on average; this ranking reaches three, giving lift~3. The constraint is five calls, not a probability threshold.
\end{workedexample}

\begin{pitfall}
Churn risk is not persuadability. Ranking metrics cannot establish that contact causes retention; use a causal design such as randomization.
\end{pitfall}

\begin{mlfigure}
\begin{tikzpicture}[baseline=(current bounding box.north)]
\begin{axis}[
  width=0.46\linewidth, height=4.5cm, axis lines=left,
  title style={font=\small\bfseries\color{MLNavy}}, title={ROC curve},
  xlabel={False positive rate}, ylabel={True positive rate},
  label style={font=\scriptsize\color{MLMuted}}, tick label style={font=\scriptsize\color{MLMuted}},
  xmin=0, xmax=1, ymin=0, ymax=1.02,
  legend style={font=\scriptsize,draw=none,fill=none,at={(0.97,0.05)},anchor=south east}
]
\addplot[MLBlue,very thick,domain=0:1,samples=60]{x^0.087}; \addlegendentry{illustrative high AUC}
\addplot[MLTeal,very thick,domain=0:1,samples=60]{x^0.389}; \addlegendentry{illustrative moderate AUC}
\addplot[MLMuted,thick,dashed,domain=0:1]{x};               \addlegendentry{random (0.50)}
\end{axis}
\end{tikzpicture}
\hfill
\begin{tikzpicture}[baseline=(current bounding box.north)]
\begin{axis}[
  width=0.46\linewidth, height=4.5cm, axis lines=left,
  title style={font=\small\bfseries\color{MLNavy}}, title={Precision-recall curve},
  xlabel={Recall}, ylabel={Precision},
  label style={font=\scriptsize\color{MLMuted}}, tick label style={font=\scriptsize\color{MLMuted}},
  xmin=0, xmax=1, ymin=0, ymax=1.02,
  legend style={font=\scriptsize,draw=none,fill=none,at={(0.03,0.05)},anchor=south west}
]
\addplot[MLBlue,very thick,domain=0.02:1,samples=60]{0.97-0.55*x^2.6}; \addlegendentry{good}
\addplot[MLTeal,very thick,domain=0.02:1,samples=60]{0.80-0.65*x^1.5}; \addlegendentry{fair}
\addplot[MLMuted,thick,dashed,domain=0:1]{0.08};                        \addlegendentry{random (= class rate)}
\end{axis}
\end{tikzpicture}
\end{mlfigure}

\begin{keyidea}
For rare positives, inspect precision--recall at the intended operating point. Strong ROC AUC can coexist with poor alert precision.
\end{keyidea}

\section{Cross-validation}
\label{sec:cross-validation}

In $k$-fold cross-validation, train on $k-1$ folds and validate on the remaining fold, rotating until all folds have been held out. Compare scores across folds.

% === VISUAL ADDITION: cross-validation-folds ===
\begin{mlfigure}
\begin{tikzpicture}[font=\scriptsize]
\node[font=\small\bfseries\color{MLNavy},anchor=west] at (0,3.55)
  {Five-fold cross-validation stays inside the development data};

\foreach \r/\yy/\lab in {0/2.85/1,1/2.33/2,2/1.81/3,3/1.29/4,4/0.77/5}{
  \node[anchor=east,text=MLMuted] at (0.85,\yy) {round \lab};
  \foreach \c/\xx in {0/1.05,1/2.43,2/3.81,3/5.19,4/6.57}{
    \pgfmathtruncatemacro{\isvalid}{\c==\r}
    \ifnum\isvalid=1
      \filldraw[fill=MLPaleOrange,draw=MLOrange,rounded corners=0.7mm]
        (\xx,\yy-0.20) rectangle (\xx+1.23,\yy+0.20);
      \node[text=MLOrange] at (\xx+0.615,\yy) {validate};
    \else
      \filldraw[fill=MLPaleBlue,draw=MLBlue!65,rounded corners=0.7mm]
        (\xx,\yy-0.20) rectangle (\xx+1.23,\yy+0.20);
      \node[text=MLBlue] at (\xx+0.615,\yy) {train};
    \fi
  }
}

\draw[MLRed,very thick,rounded corners=1mm] (8.45,0.45) rectangle (11.2,2.95);
\node[align=center,font=\small\bfseries\color{MLRed}] at (9.825,2.10)
  {Untouched\\test set};
\node[align=center,text=MLMuted] at (9.825,1.15)
  {not a sixth fold\\not used for tuning};
\draw[-{Stealth[length=2mm]},MLTeal,thick] (7.95,1.70)--(8.32,1.70);
\node[align=center,text=MLTeal] at (6.25,0.05)
  {average the five validation scores, choose the model, then evaluate the locked choice once};
\end{tikzpicture}
\end{mlfigure}
% === END VISUAL ADDITION ===

Cross-validation must preserve the data structure:

\begin{itemize}
  \item stratify folds for imbalanced classification;
  \item keep related observations in the same group fold;
  \item use forward or rolling validation for time series;
  \item fit preprocessing and feature selection separately inside each training fold.
\end{itemize}

After selecting a model and hyperparameters, refit on the allowed training data and evaluate once on the untouched test set.

\subsection{Choose the cross-validation splitter that matches the data}

\begin{mltable}
\begin{tabularx}{\linewidth}{@{}>{\bfseries\leavevmode\color{MLNavy}}p{0.27\linewidth}p{0.33\linewidth}X@{}}
\toprule
\rowcolor{MLPaleBlue!92}
\normalfont\bfseries\leavevmode\color{MLNavy}Splitter (scikit-learn) & \normalfont\bfseries\leavevmode\color{MLNavy}Use it when & \normalfont\bfseries\leavevmode\color{MLNavy}Course example \\
\midrule
\texttt{KFold} & Rows are independent and the label is a number. & Ames Housing in Lab~3 \\
\texttt{StratifiedKFold} & Classification, especially with an imbalanced label; every fold keeps the class proportions. & Titanic and breast cancer in Labs~2 and~3 \\
\texttt{GroupKFold} & Several rows belong to one person, device, or school; all of them must land in the same fold. & Repeated patient visits \\
\texttt{TimeSeriesSplit} & Rows are ordered in time; each fold trains on the past and validates on what came next. & Bike sharing in Lab~2 uses one chronological cut for the same reason \\
\bottomrule
\end{tabularx}
\end{mltable}

\begin{keyidea}
Use the same splitter for every candidate so comparisons measure model differences.
\end{keyidea}

\begin{secondpass}
The core workflow is complete. Return to probability quality, uncertainty, subgroup analysis, and shift when the application requires them.
\end{secondpass}

\section{Evaluating probabilities}

Some applications need accurate probabilities, not merely correct rankings. Binary log loss is
\[
  -\frac{1}{n}\sum_{i=1}^{n}
  [y_i\log p_i+(1-y_i)\log(1-p_i)].
\]
The Brier score is the mean squared probability error:
\[
  \operatorname{Brier}
  =\frac{1}{n}\sum_{i=1}^{n}{(p_i-y_i)}^2.
\]

\begin{workedexample}
Four predictions, with the per-row cost under each score. Log loss charges $-\log p$ when $y=1$ and $-\log(1-p)$ when $y=0$; Brier charges ${(p-y)}^2$ either way.

\begin{mltable}\small
\begin{tabular}{cccc}
\toprule
\rowcolor{MLPaleBlue!92}
True $y$ & Predicted $p$ & Log loss & Brier \\
\midrule
1 & 0.90 & $-\ln 0.90 = 0.105$ & $0.01$ \\
1 & 0.60 & $-\ln 0.60 = 0.511$ & $0.16$ \\
0 & 0.30 & $-\ln 0.70 = 0.357$ & $0.09$ \\
0 & 0.05 & $-\ln 0.95 = 0.051$ & $0.0025$ \\
\midrule
\multicolumn{2}{r}{mean} & $0.256$ & $0.0656$ \\
\bottomrule
\end{tabular}
\end{mltable}

Now add one confidently wrong row: the true answer is $1$ and the model said $0.02$. Its Brier cost is ${(0.02-1)}^2=0.960$. Its log loss cost is $-\ln 0.02 = 3.912$.

Brier is bounded by~1 per binary prediction; log loss is unbounded. Reducing the wrong prediction to $p=0.001$ changes Brier from $0.960$ to $0.998$, but log loss from $3.912$ to $6.9$.

Log loss therefore penalizes extreme confidence in a wrong answer more strongly.
\end{workedexample}

\section{Uncertainty and repeated experiments}

Metrics vary with sampled cases and training randomness. Quantify uncertainty with suitable confidence intervals, bootstrap, repeated cross-validation, or multiple seeds.

Compare models through paired errors on the same cases. Statistical significance does not imply practical value.

\section{Error and subgroup analysis}

Aggregate metrics can hide systematic failures. Inspect errors by class, label range, data source, time period, device type, location, and relevant demographic or operational groups.

Error analysis asks:

\begin{itemize}
  \item Which examples have the largest or most costly errors?
  \item Are labels ambiguous, delayed, or incorrect?
  \item Are failures concentrated in rare or underrepresented groups?
  \item Does performance degrade as data becomes older or differs from training data?
  \item Are confident mistakes more common than uncertain mistakes?
\end{itemize}

Use error analysis to identify needed changes in data, features, labels, thresholds, or models.

\section{Evaluation under distribution shift}

A test set drawn from the same source as the training data may still not look like the data the model meets after deployment. Three kinds of change have names: covariate shift (the inputs change), label shift (the class frequencies change), and concept drift (the relationship between inputs and label changes).

Stress-test missing or corrupted inputs, rare groups, seasons, and plausible worst cases. Monitor inputs, predictions, calibration, and delayed outcomes after deployment.

\begin{keyidea}
No single metric proves that a model is useful. Reliable evaluation combines a realistic split, a relevant baseline, task-specific metrics, uncertainty, subgroup analysis, and examination of actual errors.
\end{keyidea}

\begin{modelcard}{Evaluation - Practical Checklist}
\begin{tabularx}{\linewidth}{@{}>{\bfseries\leavevmode\color{MLNavy}}p{0.24\linewidth}X@{}}
Objective & Define the real decision and the cost of each kind of error. \cr
Split & Match future use; respect time, groups, location, and repeated observations. \cr
Baseline & Compare with a simple rule using the same data and metric. \cr
Uncertainty & Use cross-validation, repeated splits, or confidence intervals where appropriate. \cr
Diagnosis & Inspect errors, subgroups, calibration, threshold behavior, and distribution shift. \cr
Final report & State the data period, metric definition, uncertainty, limitations, and untouched test result. \cr
\end{tabularx}
\end{modelcard}

\begin{quickcheck}
\begin{enumerate}
  \item Describe a situation where recall should matter more than precision, and another where precision should matter more than recall.
  \item What usually happens to recall and false positives when a classification threshold is lowered?
  \item Why must the cross-validation splitter preserve groups, time order, or class balance when those structures matter to deployment?
\end{enumerate}
\end{quickcheck}
